% ===================================================================
%  TAVOLA N. 12 — La stazione. --- La ferrovia. --- Le navi.
%  Traduction 2026 d'apres texto/io, controlee sur texto/fr et
%  texto/en. Consigne : voir texto/it/10-tavola-01.tex.
%
%  LA « STAZIONE » DE CETTE PAGE EST UN RELAIS DE POSTE ROMAIN, ET
%  C'EST LA QUATRIEME FOIS QUE CETTE COLONNE PREND UN NOM ANCIEN POUR
%  UNE CHOSE NEUVE.
%
%  « Statio » etait, dans le cursus publicus de l'Empire, LE POSTE DE
%  LA ROUTE ou l'on changeait les chevaux et ou le voyageur
%  s'arretait. Le chemin de fer arrive en 1839, et l'italien ne
%  fabrique rien : il rouvre le mot. La chose est la meme a deux mille
%  ans de distance — un point sur une route ou le vehicule s'arrete et
%  ou l'homme descend — et le nom n'a pas eu a changer.
%
%  Le tableau 2 avait trouve calcio, le jeu florentin du
%  XVIe siecle donne au football anglais ; le tableau 3
%  giostra, la joute donnee au manege, et arena, le
%  sable donne au batiment. Voici le quatrieme, et il porte sur la
%  chose la plus moderne des seize planches.
%
%     LE DEPLACEMENT EST LA SIGNATURE DE CETTE COLONNE. Une langue qui
%     a trois mille ans de mots derriere elle ne forge pas et
%     n'emprunte pas : elle rouvre. Et elle rouvre d'autant plus
%     volontiers que la chose neuve ressemble davantage a une chose
%     morte.
%
%  ET LE BINARIO DE L'ALINEA 14 EST LE SEUL MOT D'EUROPE QUI
%  COMPTE LES RAILS. « Binarius », ce qui va par deux. L'allemand dit
%  « Gleis », la trace ; l'anglais « track », la piste ; le francais
%  « voie », le chemin. Tous nomment la ligne ; l'italien seul a
%  remarque qu'il y en a DEUX, et il l'a mis dans le mot.
%
%  Une exception, et elle se laisse dater par le critere du tableau 12
%  allemand. La ferrovia est un calque, non un emprunt : fer
%  plus voie. L'allemand dit « Eisenbahn », fer plus PISTE ; le
%  francais « chemin de fer », VOIE de fer. L'italien a pris « via »,
%  donc le francais et non l'allemand — et l'image, elle, n'a rien
%  d'evident : une route de metal ne se decrit pas toute seule. C'est
%  bien un calque, et il vient de Paris.
% ===================================================================

%%K t12-apar-1 apar
\VUsaut{4.0mm}
\VUtitre{15.0pt}{60}{TAVOLA N. 12}
\VUsaut{2.4mm}
\VUpk{11.4pt}{40}{La stazione. --- La ferrovia. --- Le navi.}
\VUsaut{3.4mm}
\textsc{Nota.} --- \textit{Gli orari sono diversi in ogni paese, perciò abbiamo preso
come esempio gli orari della}
\VUgras{tavola francese}.

%%K t12-01-1 p
1. --- Ho appena fatto un viaggio attraverso la Germania. Prima di partire avevo
comprato un \VUgras{orario} \textsuperscript{(15)} per sapere gli orari di partenza dei
treni. Quando ho scoperto che il treno più comodo parte da Parigi alle nove di sera e
arriva a Strasburgo all'una e tredici, mi sono preparato per il viaggio, e il giorno
dopo mi sono fatto portare alla stazione.

%%K t12-02-1 p
2. --- Quando sono entrato nel grande atrio delle partenze, dietro un vecchio
\VUgras{parroco} di campagna \textsuperscript{(2)} che portava la sua
\VUgras{borsa da viaggio} \textsuperscript{(3)}, c'erano già parecchi
\VUgras{viaggiatori} \textsuperscript{(1)}.

%%K t12-02-2 p
Siccome volevo mandare un \VUgras{telegramma} \textsuperscript{(5)}, mi sono fatto
indicare da un \VUgras{ispettore} \textsuperscript{(6)} l'\VUgras{ufficio telegrafico}
\textsuperscript{(4)}.

%%K t12-03-1 p
3. --- Prima di passare nella \VUgras{sala d'aspetto} \textsuperscript{(7)}, sono
andato a prendere un \VUgras{biglietto} di andata e ritorno \textsuperscript{(9)}.

%%K t12-04-1 p
4. --- Allo \VUgras{sportello} \textsuperscript{(8)} non ho aspettato molto, perché
c'era poca gente. Un \VUgras{soldato} \textsuperscript{(10)} che partiva in licenza mi
ha gentilmente ceduto il turno, così ho avuto il tempo di fare qualche domanda al
\VUgras{capostazione} \textsuperscript{(13)}, che stava parlando con il
\VUgras{commissario di polizia} \textsuperscript{(12)} e con il \VUgras{carabiniere}
\textsuperscript{(11)} di servizio.

%%K t12-tit-1 sub
\VUpk{10.2pt}{40}{La registrazione dei bagagli.}
\VUsaut{3.0mm}

%%K t12-05-1 p
5. --- Dopo aver comprato un giornale all'\VUgras{edicola} \textsuperscript{(14)}, ho
fatto pesare il mio \VUgras{bagaglio} \textsuperscript{(17)} --- un
\VUgras{facchino} \textsuperscript{(16)} l'aveva appena portato su su un
\VUgras{carrello} \textsuperscript{(19)}. L'\VUgras{impiegato}
\textsuperscript{(22)} addetto alla \VUgras{bilancia} \textsuperscript{(21)} mi ha
detto che il mio baule pesava trentacinque chili, il che faceva
un'\VUgras{eccedenza} di cinque chili, per la quale ho dovuto pagare un supplemento
allo \VUgras{sportello dei bagagli} \textsuperscript{(23)}.

%%K t12-06-1 p
6. --- Siccome ero in anticipo, ho avuto il tempo di guardare il movimento dei
viaggiatori nella stazione. Alcuni lasciavano o ritiravano borse al
\VUgras{deposito bagagli} \textsuperscript{(25)}; altri consultavano gli
\VUgras{orari} \textsuperscript{(26)}. I viaggiatori in arrivo correvano verso
l'\VUgras{uscita} \textsuperscript{(32)} con le loro \VUgras{valigie}
\textsuperscript{(24)} in mano. Qualcuno aspettava alla \VUgras{riconsegna bagagli}
\textsuperscript{(27)} e consegnava lo \VUgras{scontrino} \textsuperscript{(29)} per
ritirare i suoi bauli, le sue \VUgras{casse} \textsuperscript{(28)} o i suoi
\VUgras{cesti} \textsuperscript{(30)}.

%%K t12-07-1 p
7. --- I \VUgras{finanzieri} \textsuperscript{(33)} controllavano con cura i pacchi per
vedere se c'era qualcosa da dichiarare.

%%K t12-08-1 p
8. --- Qualche viaggiatore si dirigeva verso il \VUgras{buffet}
\textsuperscript{(34)}, dove un \VUgras{cameriere} \textsuperscript{(35)} li serviva in
fretta. Dovevano sbrigarsi, perché il \VUgras{treno} \textsuperscript{(36)}, un
direttissimo, non si sarebbe fermato a lungo in stazione.

%%K t12-09-1 p
9. --- Poi sono uscito sul marciapiede delle partenze e ho cercato una
\VUgras{carrozza} diretta \textsuperscript{(37)} per non dover cambiare durante il
viaggio. Sono salito su una carrozza a corridoio con uno \VUgras{scompartimento}
\textsuperscript{(38)} per fumatori e uno riservato alle «signore sole».

%%K t12-tit-2 sub
\VUpk{10.2pt}{40}{Partenza di notte.}
\VUsaut{3.0mm}

%%K t12-10-1 p
10. --- Dopo aver salito il \VUgras{predellino} \textsuperscript{(40)}, chiuso lo
\VUgras{sportello} \textsuperscript{(39)}, tirato su la \VUgras{tendina}
\textsuperscript{(41)} e messo le mie borse nella \VUgras{reticella}
\textsuperscript{(43)}, mi sono seduto sul \VUgras{sedile} \textsuperscript{(42)}.
Vedendo il \VUgras{lampista} \textsuperscript{(44)} accendere le lampade, mi sono
ricordato che dovevo passare una notte in carrozza e ho chiamato l'inserviente che
affitta i \VUgras{cuscini} \textsuperscript{(45)} per chiedergliene uno.

%%K t12-11-1 p
11. --- Il controllore è venuto a verificare i biglietti. Gli ho chiesto perché non
eravamo partiti, dato che era già ora. Mi ha risposto che il \VUgras{capotreno}
\textsuperscript{(46)} era occupato a far caricare delle biciclette nel
\VUgras{bagagliaio} \textsuperscript{(47)}. Inoltre non erano ancora stati cambiati
tutti gli \VUgras{scaldini} \textsuperscript{(48)}.

%%K t12-12-1 p
12. --- Ma stavamo per partire: il \VUgras{fuochista} \textsuperscript{(50)}, in cima
al suo \VUgras{tender} \textsuperscript{(49)}, prendeva carbone per alimentare il fuoco
della \VUgras{locomotiva} \textsuperscript{(54)}, e \VUgras{il macchinista}
\textsuperscript{(51)} aspettava solo il segnale del \VUgras{vicecapostazione}
\textsuperscript{(52)}, che stava sul \VUgras{marciapiede} \textsuperscript{(53)}.

%%K t12-13-1 p
13. --- La \VUgras{caldaia} \textsuperscript{(56)} della locomotiva rombava sordamente;
il \VUgras{fumaiolo} \textsuperscript{(55)} buttava fuori torrenti di fumo misto a
\VUgras{vapore} \textsuperscript{(60)}. Un \VUgras{fischio} \textsuperscript{(59)}
acuto è risuonato e gli \VUgras{stantuffi} \textsuperscript{(58)} si sono messi in
moto, facendo girare le \VUgras{ruote} \textsuperscript{(57)} della pesante macchina.

%%K t12-tit-3 sub
\VUsaut{3.0mm}
\VUpk{10.2pt}{40}{Si parte!}
\VUsaut{3.0mm}

%%K t12-14-1 p
14. --- Il treno, che correva sulle \VUgras{rotaie} \textsuperscript{(64)}, prendeva
velocità; i \VUgras{marciapiedi} \textsuperscript{(65)} finivano; siamo passati
davanti alle \VUgras{ritirate} \textsuperscript{(66)}, e anche davanti a una fila di
carrozze ferme su un altro \VUgras{binario} \textsuperscript{(63)}: \VUgras{vagoni
letto} \textsuperscript{(67)}, un \VUgras{vagone ristorante} \textsuperscript{(68)},
un \VUgras{vagone postale} \textsuperscript{(69)}.

%%K t12-noto-1 noto
\VUnotes{143.2mm}{%
\textit{(*) Nota.} --- \textit{Il viaggio è naturalmente descritto com'era lungo il
confine di prima della guerra.}%
}

%%K t12-15-1 p
15. --- Poi ci è passato davanti lo \VUgras{scalo merci} \textsuperscript{(71)}. Sotto
le tettoie gli impiegati caricavano le \VUgras{merci} \textsuperscript{(72)} spedite a
piccola velocità. Alcuni \VUgras{manovratori} \textsuperscript{(76)} vicino al
\VUgras{serbatoio dell'acqua} \textsuperscript{(73)} spostavano \VUgras{carri
bestiame} \textsuperscript{(75)} e \VUgras{pianali} \textsuperscript{(79)} per formare
un \VUgras{treno merci} \textsuperscript{(80)}.

%%K t12-16-1 p
16. --- Abbiamo superato l'ultimo \VUgras{scambio} \textsuperscript{(78)}, accanto al
quale stava il deviatore; siamo passati davanti al \VUgras{segnale}
\textsuperscript{(86)} e la stazione è scomparsa in lontananza.

%%K t12-17-1 p
17. --- Poco dopo siamo usciti sul \VUgras{ponte di ferro} \textsuperscript{(74)} che
attraversa il \VUgras{fiume} \textsuperscript{(81)}. Passato il ponte, siamo andati
lungo il fiume, sul quale potenti \VUgras{rimorchiatori} \textsuperscript{(82)}
tiravano pesanti \VUgras{chiatte} \textsuperscript{(83)}. Abbiamo visto anche un
\VUgras{battello passeggeri} \textsuperscript{(84)} che accostava a un
\VUgras{pontile} \textsuperscript{(85)}.

%%K t12-18-1 p
18. --- Dopo aver attraversato a tutta velocità parecchie stazioni, siamo passati per
qualche \VUgras{galleria} \textsuperscript{(87)} e ci siamo lasciati dietro
innumerevoli \VUgras{casette} \textsuperscript{(88)} dei \VUgras{casellanti}.
Cullato dal dondolio del treno, sono caduto in un sonno profondo.

%%K t12-tit-4 sub
\VUsaut{3.0mm}
\VUpk{10.2pt}{40}{Stazione di Deutsch-Avricourt.}
\VUsaut{3.0mm}

%%K t12-19-1 p
19. --- Mi ha svegliato di colpo la voce di un impiegato che gridava:
«Deutsch-Avricourt! \textsuperscript{(*)} Si cambia!» Era il controllo doganale con
tutte le formalità noiose della frontiera. Ho dovuto mettere l'orologio sull'ora
dell'\VUgras{orologio} \textsuperscript{(89)} della stazione, che era avanti di
cinquantacinque minuti; mezzo addormentato, distinguevo a fatica la \VUgras{lancetta
grande} \textsuperscript{(90)} dalla \VUgras{piccola} \textsuperscript{(91)}; il
\VUgras{quadrante} \textsuperscript{(92)} stesso era tutto offuscato dalla nebbia del
mattino.

%%K t12-20-1 p
20. --- Mentre i doganieri facevano il loro controllo, ho guardato sbadigliando i
\VUgras{manifesti} \textsuperscript{(93)} che ornano i muri della stazione. Ho fatto
colazione per passare il tempo, ho consultato la carta della \VUgras{rete}
\textsuperscript{(94)} tedesca per vedere quanto mancava a Strasburgo, e sono risalito
nella mia carrozza, rincuorato dalla speranza di arrivare presto a destinazione.
\VUsaut{6.0mm}
\VUfilet{22mm}
